\Article{37}
In Art.\ 15.\ 16 haben wir eine Methode erläutert, um die Genauigkeit der Beobachtungen näherungsweise zu bestimmen\footnote{Die Untersuchung desselben Gegenstands, die wir in einer früheren Abhandlung (\textit{Bestimmung der Genauigkeit der Beobachtungen. Zeitschrift für Astronomie und verwandte Wissenschaften} Vol.\ I, p.\ 185) vorgelegt hatten, beruhte auf derselben Hypothese über die Beschaffenheit der Funktion, welche die Wahrscheinlichkeit der Fehler ausdrückt, auf die wir in der Theoria motus corporum coelestium die Methode der kleinsten Quadrate gegründet hatten (siehe Art.\ 9, III).}. Diese Methode setzt jedoch voraus, dass eine hinreichend große Zahl der tatsächlich aufgetretenen Fehler genau bekannt ist; diese Bedingung wird, streng genommen, äußerst selten, um nicht zu sagen niemals, erfüllt sein. Hängen nämlich die Größen, deren Näherungswerte durch Beobachtungen bekannt geworden sind, nach einem bekannten Gesetz von einer oder mehreren unbekannten Größen ab, so lassen sich deren plausibelste Werte mit der Methode der kleinsten Quadrate ermitteln. Von den daraus berechneten Werten der Größen, die Gegenstand der Beobachtungen waren, wird dann angenommen, dass sie nur sehr wenig von den
\SourceLeaf{052}%
wahren Werten abweichen. Ihre Unterschiede gegenüber den beobachteten Werten darf man daher mit umso größerem Recht als die wahren Beobachtungsfehler annehmen, je größer die Zahl der Beobachtungen ist. Diesem Verfahren sind alle Rechner gefolgt, die es unternommen haben, die Genauigkeit der Beobachtungen in konkreten Fällen a posteriori zu schätzen: Es ist jedoch offenbar theoretisch falsch; und obwohl es in vielen Fällen für praktische Zwecke ausreichen mag, kann es in anderen zu gewaltigen Fehlern führen. Dieser Gegenstand verdient daher in höchstem Maße eine genauere Untersuchung.

Wir behalten in dieser Untersuchung die seit Art.\ 19 verwendeten Bezeichnungen bei. Das besprochene Verfahren betrachtet die Größen \(A,B,C\) etc. als die wahren Werte von \(x,y,z\) und folglich \(\lambda,\lambda',\lambda''\) etc. als die wahren Werte der Funktionen \(v,v',v''\) etc. Haben alle Beobachtungen dieselbe Genauigkeit und wird ihr Gewicht \(p=p'=p''\) etc. als Einheit genommen, so stellen dieselben Größen bei umgekehrten Vorzeichen unter jener Annahme die Beobachtungsfehler dar. Die Vorschriften von Art.\ 15 liefern dann für den mittleren Beobachtungsfehler \(m\)
\[
=\BareRoot{\frac{\lambda\lambda+\lambda'\lambda'+\lambda''\lambda''+\text{etc.}}{\pi}}=\BareRoot{\frac{M}{\pi}}
\]
Ist die Genauigkeit der Beobachtungen nicht dieselbe, so stellen die Größen \(-\lambda,-\lambda',-\lambda''\) etc. die mit den Quadratwurzeln der Gewichte multiplizierten Beobachtungsfehler dar; die Vorschriften von Art.\ 16 würden auf dieselbe Formel \(\BareRoot{\frac{M}{\pi}}\) führen, die nun den mittleren Fehler solcher Beobachtungen bezeichnet, denen das Gewicht \(=1\) zugewiesen wird. Offenbar würde aber eine genaue Rechnung erfordern, dass man anstelle der Größen \(\lambda,\lambda',\lambda''\) etc. die Werte der Funktionen \(v,v',v''\) etc. verwendet, die sich aus den wahren Werten von \(x,y,z\) etc. ergeben, d.\ h. anstelle von \(M\) den Wert der Funktion \(\Omega\), der den wahren Werten von \(x,y,z\) etc. entspricht. Dieser lässt sich zwar nicht angeben; dennoch sind wir sicher, dass er größer ist als \(M\) (das ja der kleinstmögliche Wert ist), ausgenommen den unendlich unwahrscheinlichen Fall, dass die plausibelsten Werte der Unbekannten genau mit den wahren übereinstimmen. Im Allgemeinen können wir daher behaupten, dass das übliche Verfahren einen zu kleinen mittleren Fehler liefert oder den Beobachtungen eine zu große Genauigkeit zuschreibt. Sehen wir nun, was die strenge Theorie lehrt.

\Article{38}
Vor allem ist zu untersuchen, wie \(M\) von den wahren Beobachtungsfehlern abhängt. Bezeichnen wir diese wie in Art.\ 28 mit \(e,e',e''\) etc. und setzen wir zur größeren Einfachheit
\SourceLeaf{053}
\[
\begin{gathered}
e\BareRoot{p}=\varepsilon,\quad e'\BareRoot{p'}=\varepsilon',\quad e''\BareRoot{p''}=\varepsilon''\quad\text{etc., sowie}\\
m\BareRoot{p}=m'\BareRoot{p'}=m''\BareRoot{p''}\quad\text{etc.}=\mu
\end{gathered}
\]
Ferner seien die wahren Werte von \(x,y,z\) etc. jeweils \(A-x^0,\allowbreak B-y^0,\allowbreak C-z^0\) etc.; ihnen mögen für \(\xi,\eta,\zeta\) etc. die Werte \(-\xi^0,-\eta^0,-\zeta^0\) etc. entsprechen. Offenbar entsprechen denselben Werten für \(v,v',v''\) etc. die Werte \(-\varepsilon,-\varepsilon',-\varepsilon''\) etc., sodass gilt
\[
\begin{aligned}
\xi^0&=a\varepsilon+a'\varepsilon'+a''\varepsilon''+\text{etc.}\\
\eta^0&=b\varepsilon+b'\varepsilon'+b''\varepsilon''+\text{etc.}\\
\zeta^0&=c\varepsilon+c'\varepsilon'+c''\varepsilon''+\text{etc.}
\end{aligned}
\]
etc. sowie
\[
\begin{aligned}
x^0&=\alpha\varepsilon+\alpha'\varepsilon'+\alpha''\varepsilon''+\text{etc.}\\
y^0&=\beta\varepsilon+\beta'\varepsilon'+\beta''\varepsilon''+\text{etc.}\\
z^0&=\gamma\varepsilon+\gamma'\varepsilon'+\gamma''\varepsilon''+\text{etc.}
\end{aligned}
\]
Setzen wir schließlich
\[
\Omega^0=\varepsilon\varepsilon+\varepsilon'\varepsilon'+\varepsilon''\varepsilon''+\text{etc.}
\]
sodass \(\Omega^0\) gleich dem Wert der Funktion \(\Omega\) ist, der den wahren Werten von \(x,y,z\) etc. entspricht. Da unbestimmt gilt
\[
\Omega=M+(x-A)\xi+(y-B)\eta+(z-C)\zeta+\text{etc.}
\]
gilt auch
\[
M=\Omega^0-x^0\xi^0-y^0\eta^0-z^0\zeta^0-\text{etc.}
\]
Daraus ist klar, dass \(M\) nach der Entwicklung eine homogene Funktion zweiter Ordnung der Fehler \(e,e',e''\) etc. ist, die je nach den verschiedenen Werten der Fehler größer oder kleiner werden kann. Solange uns aber die Größe der Fehler unbekannt bleibt, wird es zweckmäßig sein, diese Funktion unbestimmt zu betrachten und vor allem nach den Grundsätzen der Wahrscheinlichkeitsrechnung ihren Mittelwert anzugeben. Diesen finden wir, wenn wir anstelle der Quadrate \(ee,e'e',e''e''\) etc. jeweils \(mm,m'm',m''m''\) etc. schreiben, die Produkte \(ee',ee'',e'e''\) etc. dagegen gänzlich weglassen; oder, was dasselbe ist, wenn wir anstelle jedes Quadrats \(\varepsilon\varepsilon,\varepsilon'\varepsilon',\varepsilon''\varepsilon''\) etc. den Wert \(\mu\mu\) schreiben und die Produkte \(\varepsilon\varepsilon',\varepsilon\varepsilon'',\varepsilon'\varepsilon''\) etc. vollständig außer Acht lassen. Auf diese Weise ergibt sich aus dem Glied \(\Omega^0\) offenbar \(\pi\mu\mu\); das Glied \(-x^0\xi^0\) liefert
\[
-(a\alpha+a'\alpha'+a''\alpha''+\text{etc.})\mu\mu=-\mu\mu
\]
\SourceLeaf{054}%
und ebenso liefert jeder der übrigen Teile \(-\mu\mu\), sodass der gesamte Mittelwert \(=(\pi-\rho)\mu\mu\) wird, wobei \(\pi\) die Zahl der Beobachtungen und \(\rho\) die Zahl der Unbekannten bezeichnet. Der wahre Wert von \(M\) kann freilich, je nachdem, welche Fehler der Zufall hervorgebracht hat, größer oder kleiner als der Mittelwert ausfallen. Die Abweichung wird jedoch umso weniger ins Gewicht fallen, je größer die Zahl der Beobachtungen ist, sodass man als Näherungswert von \(\mu\) annehmen darf
\[
\BareRoot{\frac{M}{\pi-\rho}}
\]
Der Wert von \(\mu\), der sich aus dem im vorigen Artikel besprochenen fehlerhaften Verfahren ergibt, muss daher im Verhältnis der Größe \(\BareRoot{(\pi-\rho)}\) zu \(\BareRoot{\pi}\) vergrößert werden.

\clearpage
